\documentclass[11pt,a4paper]{article}

% ---------- Core encoding / fonts ----------
\usepackage[utf8]{inputenc}
\usepackage[T1]{fontenc}
\usepackage{lmodern}
\usepackage{microtype}
\usepackage[strings]{underscore}   % literal underscores in text mode, no manual escaping

% ---------- Layout ----------
% Generous top/bottom margins so the header/footer sit clearly between
% the page-edge border (drawn 1cm in, see below) and the body text --
% no overlap in either direction.
\usepackage[a4paper,left=2.6cm,right=2.6cm,top=3.4cm,bottom=3.2cm,
            headheight=14pt,headsep=22pt,footskip=28pt]{geometry}
\usepackage{tikz}
\usetikzlibrary{calc}
\usepackage{eso-pic}
\usepackage{fancyhdr}
\usepackage{titlesec}
\usepackage{enumitem}
\usepackage{listings}
\usepackage{xcolor}
\usepackage{hyperref}
\usepackage{parskip}
\usepackage{seqsplit}

% ---------- Strictly black & white ----------
\hypersetup{
  colorlinks=true,
  linkcolor=black,
  citecolor=black,
  urlcolor=black,
  pdftitle={Kernel Borderlands -- Bug Report},
  pdfauthor={Kernel Borderlands Team}
}

% ---------- Four-side page border on every page ----------
\newcommand{\pageborder}{%
  \begin{tikzpicture}[remember picture, overlay]
    \draw[black, line width=0.7pt]
      ($(current page.north west)+(1.0cm,-1.0cm)$) rectangle
      ($(current page.south east)+(-1.0cm,1.0cm)$);
  \end{tikzpicture}%
}
\AddToShipoutPictureBG{\pageborder}

% ---------- Headers / footers ----------
\pagestyle{fancy}
\fancyhf{}
\renewcommand{\headrulewidth}{0.4pt}
\renewcommand{\footrulewidth}{0pt}
\fancyhead[L]{\small\scshape Kernel Borderlands}
\fancyhead[R]{\small\scshape Bug Report}
\fancyfoot[C]{\small\thepage}

\fancypagestyle{plain}{%
  \fancyhf{}
  \renewcommand{\headrulewidth}{0pt}
  \fancyfoot[C]{\small\thepage}
}

% ---------- Section styling (clean, sharp, no color) ----------
\titleformat{\section}
  {\normalfont\Large\bfseries}{\thesection}{0.8em}{}
\titleformat{\subsection}
  {\normalfont\large\bfseries}{}{0em}{}
\titlespacing*{\subsection}{0pt}{1.4em}{0.6em}

% ---------- Code listings: plain monospace, no color ----------
\lstset{
  basicstyle=\ttfamily\small,
  breaklines=true,
  breakatwhitespace=false,
  columns=fullflexible,
  frame=single,
  framerule=0.4pt,
  rulecolor=\color{black},
  backgroundcolor=\color{white},
  showstringspaces=false,
  xleftmargin=1pt,
  xrightmargin=1pt,
  tabsize=2,
  aboveskip=0.8em,
  belowskip=0.8em
}

% ---------- Helper macros for bug entries ----------
\newcommand{\bugmeta}[3]{%
  \noindent\textbf{Severity:} #1 \\
  \textbf{Component:} \ttfamily\seqsplit{#2}\normalfont \\
  \textbf{Status:} #3
  \par\vspace{0.5em}
}

\newcommand{\entrysep}{\vspace{0.4em}\noindent\rule{\linewidth}{0.3pt}\vspace{1.1em}}

% A printed name with a styled, greyed-out "signature" underneath --
% uses the standard Zapf Chancery PSNFSS font for a script/signature look,
% no external font install needed.
% #1 = name, #2 = roll number
\newcommand{\devcredit}[2]{%
  \begin{minipage}[t]{3.2cm}
    \centering
    {\small #1}\\[0.35em]
    {\color{gray!65}\fontfamily{pzc}\selectfont\large #2}
  \end{minipage}%
}

\setlist[itemize]{leftmargin=1.4em, itemsep=0.2em}
\setlist[enumerate]{leftmargin=1.6em, itemsep=0.2em}

\begin{document}

% ================= TITLE PAGE =================
\begin{titlepage}
\thispagestyle{empty}
\centering
\vspace*{\fill}

{\small\textls[280]{KERNEL BORDERLANDS}}

\vspace{1.6cm}

{\fontsize{40}{48}\selectfont\bfseries Bug Report}

\vspace{0.8em}

{\normalsize A Source-Verified Audit of Open, Resolved, and Historical Defects}

\vspace{1.2cm}

{\small Compiled for August 2nd, 2026}

\vspace*{\fill}

\begin{minipage}{0.8\linewidth}
\centering
\small
\textbf{Scope} \\[0.3em]
Direct source audit across all four subsystems --- \texttt{kb-core},
\texttt{kb-aads}, \texttt{kb-control-plane}, \texttt{kb-op} --- cross-checked
against \seqsplit{\texttt{docs/development/core-control/control-plane-catalog.md}} and
\texttt{CHANGELOG.md}. Every finding was independently verified by reading
the current source directly, not taken on faith from prior documentation.
\end{minipage}

\vspace{1.4cm}

{\small \scshape Kernel Borderlands Team}\\[0.9em]
\devcredit{Karthik}{2311CS040080}\hspace{0.6cm}\devcredit{PardhuVarma}{2311CS040089}\hspace{0.6cm}\devcredit{Tejaswini}{2311CS040077}\hspace{0.6cm}\devcredit{Rupa}{2311CS040082}

\vspace{1.5cm}
\end{titlepage}

% ================= TOC =================
\tableofcontents
\thispagestyle{plain}
\clearpage

% ================= INTRODUCTION =================
\section*{Introduction}
\addcontentsline{toc}{section}{Introduction}

This report covers real, code-level defects and security gaps in the
Kernel Borderlands product itself --- sourced from
\seqsplit{\texttt{docs/development/core-control/control-plane-catalog.md}} (the
control plane's own severity-graded \emph{Open Gaps} register) and
\texttt{CHANGELOG.md}, cross-checked directly against the current source
to confirm what is actually still open versus already fixed since those
documents were written.

Dev-environment and tooling issues encountered during development are
intentionally excluded from this report --- see
\seqsplit{\texttt{docs/reports/kb-aads/containment-rl-2026-08-02.md}} for those. This
revision adds a direct source audit across all four subsystems
(\texttt{kb-core}, \texttt{kb-aads}, \texttt{kb-control-plane},
\texttt{kb-op}); every finding below was independently verified by reading
the actual current source, not taken on faith from any prior document.

\clearpage

% ================= OPEN BUGS =================
\section{Open --- Confirmed Still Present in Current Source}

\subsection*{BUG-001 --- Unauthenticated, Wildcard-CORS HTTP API Can Terminate/Restore Containment on Any PID}
\addcontentsline{toc}{subsection}{BUG-001 --- Unauthenticated HTTP API}
\bugmeta{Critical (Security)}{kb-control-plane/internal/controlplane/\{controlplane.go, http.go\}, triggered from kb-op/kb-dashboard/src/App.tsx}{Open}

\textbf{Description.} \texttt{kbd} starts an HTTP API on \texttt{:8080} ---
bound to \textbf{all interfaces}, not loopback-only --- serving
\texttt{/api/isolate} and \texttt{/api/restore} with \textbf{zero
authentication}, and a CORS handler that sets
\texttt{Access-Control-Allow-Origin: *} on every response.

\textbf{Confirmed in source:}
\begin{lstlisting}[language=Go]
// controlplane.go:184
if err := cp.StartHTTPServer(":8080"); err != nil {
    // all interfaces, not 127.0.0.1:8080

// http.go:46
w.Header().Set("Access-Control-Allow-Origin", "*")
    // every origin, no check

// http.go:398-423, handleIsolate -- no auth check anywhere
var req struct{ Pid uint32 `json:"pid"` }
json.NewDecoder(r.Body).Decode(&req)
s.cp.enforcer.Contain(req.Pid, uint32(pb.ContainmentLevel_TERMINATE),
    "Manual isolation from dashboard")
\end{lstlisting}
\texttt{handleRestore} (same file, near line 425) is identical in shape,
calling \seqsplit{\texttt{Contain(..., ContainmentLevel\_NONE, ...)}}.

\textbf{Impact.} Any actor with network reach to port 8080 --- or, because
of the wildcard CORS header, \emph{any webpage a browser loads at all},
from any origin --- can POST \texttt{\{"pid": N\}} to \texttt{/api/isolate}
and terminate containment on an arbitrary PID with no credential, token, or
origin check whatsoever. This is strictly worse than BUG-011/BUG-012
(which at least require reaching a Unix socket): this is an open TCP
listener with explicit cross-origin permission, reachable from anywhere
the port is routable. Combined with BUG-011 (no protected-PID gate on the
Go side), this path can also be pointed at PID 1 or \texttt{kbd}'s own PID
with no additional resistance. Not mentioned anywhere in
\texttt{control-plane-catalog.md}, which only discusses this daemon's
gRPC-facing gaps --- its HTTP API's auth posture appears to have gone
unreviewed entirely.

\textbf{Suggested fix.} Bind to \texttt{127.0.0.1:8080} unless explicit
remote-dashboard access is configured; require a real auth token/session
on \texttt{/api/isolate} and \texttt{/api/restore} (and any other
state-mutating route); replace the wildcard CORS header with an explicit
allowed-origin list.

\entrysep

\subsection*{BUG-002 --- UpsertProcessState Silently Resets an Actively-Contained Process Back to NONE}
\addcontentsline{toc}{subsection}{BUG-002 --- Containment reset on telemetry update}
\bugmeta{Critical}{kb-control-plane/internal/store/process.go}{Open}

\textbf{Description.} \texttt{UpsertProcessState} --- called on the hot
path, continuously, for every incoming telemetry update from
\texttt{kb-core} --- builds a \textbf{brand-new} \texttt{CachedState\{...\}}
literal and overwrites L1 wholesale. Unlike every other mutator in this
file (\texttt{SetContainment}, \texttt{SetZone},
\texttt{InsertZoneTransition}, all of which correctly load-then-mutate-one-
field), this one omits the \texttt{Containment} field entirely, so it is
always the Go zero-value (\texttt{0} = \texttt{NONE}) after any telemetry
update.

\textbf{Confirmed in source:}
\begin{lstlisting}[language=Go]
// SetContainment (correct pattern):
if v, ok := s.l1.Load(pid); ok {
    updated := *v.(*CachedState)
    updated.Containment = level
    s.l1.Store(pid, &updated)
}

// UpsertProcessState (the bug):
s.l1.Store(msg.PID, &CachedState{
    PID: msg.PID, PPID: msg.PPID, UID: msg.UID, Comm: msg.Comm,
    StartTimeNs: msg.StartTimeNs, LastUpdatedNs: msg.LastUpdatedNs,
    DimScore: msg.DimScore, CompositeScore: msg.CompositeScore,
    EMAScore: msg.EMAScore,
    SyscallEntropyLifetime: msg.SyscallEntropyLifetime,
    Zone: msg.Zone, EventCount: msg.EventCount,
    // Containment: not set -- zero value, i.e. NONE
})
\end{lstlisting}

\textbf{Impact.} An operator (or agent) calls
\texttt{SetContainment(pid, TERMINATE)} $\rightarrow$ correctly recorded.
\texttt{kb-core} keeps sending routine telemetry for that PID (containment
does not stop monitoring) $\rightarrow$ the very next \texttt{ProcessState}
update overwrites L1 with \texttt{Containment: NONE}. Every reader of
\texttt{GetProcessState} --- \texttt{kb-tui}, \texttt{kbctl}, the dashboard,
\texttt{kb-aads} --- now reports an actively-contained process as
uncontained, within one telemetry cycle (typically sub-second). This is a
state-consistency bug with direct operational-security impact: the
system's own status reporting cannot be trusted for containment state.

\textbf{Suggested fix.} Load the existing \texttt{CachedState} and copy
forward \texttt{Containment} (and ideally do a true partial-field update,
not full replacement) the same way \texttt{SetContainment}/\texttt{SetZone}
already do.

\entrysep

\subsection*{BUG-003 --- JJE Consensus: Jury's Containment Vote Can Never Trigger on Any Real Input}
\addcontentsline{toc}{subsection}{BUG-003 --- Jury threshold scale mismatch}
\bugmeta{High}{kb-aads/consensus/jje.py, JuryAgent.evaluate\_and\_vote}{Open}

\textbf{Description.} \texttt{vote = "CONTAIN" if score > 75.0 else
"ALLOW"}, where \texttt{score = alert\_payload.get("confidence", 0.0)}.
Every other confidence value in this system is a $0.0$--$1.0$ float ---
confirmed against \seqsplit{\texttt{kb-control-plane/internal/controlplane/grpc.go}}'s
\texttt{SubmitAgentDecision} (\texttt{"confidence \%.2f below 0.85
threshold"}) and \texttt{kb.proto}'s confidence fields. A real alert with
confidence $0.95$ (about as certain as this system ever gets) is not
$> 75.0$ --- so under the scale used everywhere else, this condition is
never true for any realistic payload.

\textbf{Impact.} The Judge $\rightarrow$ Jury $\rightarrow$ Executor
consensus containment path --- the mechanism this project's own
architecture treats as the primary quorum-gated containment decision --- is
effectively dead code. Jury never votes CONTAIN, so
\texttt{coordinate\_consensus}'s \texttt{contain\_votes / total\_votes >
0.5} check can never pass, so \texttt{execute\_quarantine} is never
triggered via this path. The identical threshold also appears in
\seqsplit{\texttt{docs/development/control-aads/aads-development-plan.md:300}},
suggesting it was copied from an early design sketch and never reconciled
with the 0--1 convention the rest of the system settled on.

\textbf{Suggested fix.} Change the threshold to the 0--1 scale (e.g.
\texttt{score > 0.75}), and add a test asserting a realistic
high-confidence payload (\texttt{confidence: 0.9}+) actually produces a
CONTAIN vote --- the kind of test that would have caught this immediately.

\entrysep

\subsection*{BUG-004 --- kb\_lsm\_file\_open Fails Open on Path-Resolution Failure, Before the Containment-Level Check Ever Runs}
\addcontentsline{toc}{subsection}{BUG-004 --- LSM fail-open on bpf\_d\_path failure}
\bugmeta{High}{kb-core/ebpf/kbd\_sensor.bpf.c}{Open}

\textbf{Description.}
\begin{lstlisting}[language=C]
int len = bpf_d_path(&file->f_path, path_buf, sizeof(path_buf));
if (len < 0) return 0;   // allow -- before containment level is checked
...
if (level) {
    if (*level >= 3) {
        return -13; // Block ALL file open requests (full sandbox quarantine)
    }
    ...
\end{lstlisting}
If \texttt{bpf\_d\_path()} fails (a real, documented failure mode for
certain special files, pseudo-filesystems, or under kernel memory pressure
--- not hypothetical), the function returns "allow" immediately, before the
containment-level $\geq 3$ "block all file opens" check ever executes.

\textbf{Impact.} A process placed at containment level 3 or 4 --- which the
code's own comment states should have "ALL file open requests" blocked for
full sandbox quarantine --- has that guarantee silently bypassed for any
file whose path resolution triggers this failure mode. This is distinct
from the already-fixed zero-padding bug (HIST-3, which was about
sensitive-path \emph{string matching}); this is about the level-3+ block
never running at all in this code path.

\textbf{Suggested fix.} On \texttt{bpf\_d\_path()} failure, fail closed for
any contained process (check \texttt{level} first, before attempting path
resolution, and block/return \texttt{-13} for level $\geq 3$ regardless of
whether the path was resolved) rather than allowing by default.

\entrysep

\subsection*{BUG-005 --- VerifyChain() Discards the SQL Scan Error, Risking a False Tamper Verdict}
\addcontentsline{toc}{subsection}{BUG-005 --- VerifyChain ignores Scan error}
\bugmeta{Medium-High}{kb-control-plane/internal/audit/audit.go}{Open}

\textbf{Description.}
\begin{lstlisting}[language=Go]
rows.Scan(&ts, &action, &subject, &actor, &reason, &prevH, &entryH)
\end{lstlisting}
The returned error is never checked. On a scan failure (unexpected NULL,
type mismatch, a corrupted row), the loop proceeds with zero-valued
fields, hashes garbage, and reports \texttt{"chain broken at entry N"} ---
a false tamper alert --- rather than the actual I/O/schema error. This
directly undermines RESOLVED-1: the verifier that was just wired up has a
latent bug that can produce an incorrect verdict.

\textbf{Suggested fix.} \texttt{if err := rows.Scan(...); err != nil \{
return false, count, err \}}.

\entrysep

\subsection*{BUG-006 --- kb-aads Jury Pool Size Is Hardcoded, Silently Ignoring Its Own Config File}
\addcontentsline{toc}{subsection}{BUG-006 --- Jury pool size ignores config}
\bugmeta{Medium}{kb-aads/consensus/jje.py, JudgeAgent.coordinate\_consensus}{Open}

\textbf{Description.} \texttt{config/agents.yaml} declares \texttt{jury:
pool\_size: 5}, but \texttt{coordinate\_consensus} hardcodes
\texttt{jury\_pool = [JuryAgent.remote(f"jury-\{i\}") for i in range(5)]}
and never reads that config value. Currently invisible only because the
numbers happen to match by coincidence.

\textbf{Impact.} An operator changing \texttt{pool\_size} in config,
expecting it to take effect, silently gets no change at all.

\textbf{Suggested fix.} Read \texttt{config/agents.yaml}'s
\texttt{jury.pool\_size} instead of the literal \texttt{5}.

\entrysep

\subsection*{BUG-007 --- Consensus Round Has No Fault Tolerance for a Crashed Juror}
\addcontentsline{toc}{subsection}{BUG-007 --- No fault tolerance in consensus round}
\bugmeta{Medium}{kb-aads/consensus/jje.py, JudgeAgent.coordinate\_consensus}{Open}

\textbf{Description.} \texttt{votes = ray.get(vote\_futures)} has no error
handling. If any spawned \texttt{JuryAgent} dies or its \texttt{.remote()}
call raises, the exception propagates uncaught and kills the entire
consensus round.

\textbf{Impact.} One crashed/misbehaving juror blocks the whole
containment decision for that alert, rather than the round degrading
gracefully (excluding the failed juror, re-normalizing remaining weights).

\textbf{Suggested fix.} Wrap in a per-future try/except (or
\texttt{ray.wait} with a timeout) and proceed with whichever jurors
actually responded.

\entrysep

\subsection*{BUG-008 --- Unbounded Per-PID eBPF Maps Are Never Cleaned Up, Unlike Their CPM/CWP Siblings}
\addcontentsline{toc}{subsection}{BUG-008 --- eBPF map leak}
\bugmeta{Medium}{kb-core/ebpf/kbd\_sensor.bpf.c}{Open}

\textbf{Description.} \texttt{kb\_handle\_exit} explicitly deletes
\texttt{protected\_pids\_map} and \texttt{protected\_workloads\_map}
entries on process exit (correctly, for PID-reuse safety) --- but three
other PID-keyed maps are never touched by any exit path:
\texttt{kb\_syscall\_totals} (10240 entries), \texttt{kb\_syscall\_counts}
(655360 entries), \texttt{kb\_cred\_prev} (10240 entries). All
\texttt{bpf\_map\_update\_elem} calls against them also ignore their
return value.

\textbf{Impact.} On a long-running host with normal PID churn, these maps
fill permanently. Once full, new-PID inserts silently fail (\texttt{-E2BIG},
unchecked) --- syscall-volume telemetry and the privilege-escalation
UID-diffing in \texttt{kb\_handle\_commit\_creds} stop updating for any new
process from that point on, with no error surfaced anywhere.

\textbf{Suggested fix.} Delete corresponding entries from all three maps
in \texttt{kb\_handle\_exit}, matching the pattern already used for
\texttt{protected\_pids\_map}/\texttt{protected\_workloads\_map}.

\entrysep

\subsection*{BUG-009 --- Dashboard Hardcodes http://localhost:8080, Compounding BUG-001 if Ever Pointed at a Real Host}
\addcontentsline{toc}{subsection}{BUG-009 --- Dashboard hardcoded localhost}
\bugmeta{Low (Medium once BUG-001 is fixed and this is deployed remotely)}{kb-op/kb-dashboard/src/App.tsx (7 call sites)}{Open}

\textbf{Description.} The dashboard's API base URL is hardcoded to
\texttt{localhost:8080} rather than being environment-configurable.

\textbf{Impact.} Breaks the dashboard for any non-localhost deployment
today. More importantly, it means the natural fix anyone reaches for
("just change the hardcoded URL to the real host") makes BUG-001's
vulnerability remotely exploitable instead of LAN-local, unless BUG-001 is
fixed first.

\textbf{Suggested fix.} Externalize to a build-time/runtime environment
variable.

\entrysep

\subsection*{BUG-010 --- AgentDecision.AuthorizedBy Is Captured but Never Validated}
\addcontentsline{toc}{subsection}{BUG-010 --- AuthorizedBy unvalidated}
\bugmeta{High (Security)}{kb-control-plane/internal/controlplane/grpc.go, SubmitAgentDecision}{Open}

\textbf{Description.} \texttt{AuthorizedBy} is logged into the audit trail
as if it were a meaningful authorization claim, but it is just a
caller-supplied string field on the incoming \texttt{AgentDecision}
message --- nothing checks it against an allowlist, a signature, or any
real authority before acting on it.

\textbf{Confirmed in source} (\texttt{grpc.go}):
\begin{lstlisting}[language=Go]
cp.audit.Log(
    fmt.Sprintf("AGENT_%s", d.Action),
    fmt.Sprintf("pid=%d agent=%s conf=%.2f auth=%v",
        d.Pid, d.AgentId, d.Confidence, d.AuthorizedBy),
    d.AgentId, "",
)
\end{lstlisting}
\texttt{d.AuthorizedBy} is read into the log string and nowhere else.

\textbf{Impact.} Combined with BUG-011, any process able to reach
\texttt{kba.sock} can submit a fabricated
\texttt{TERMINATE}/\texttt{NAMESPACE}/\texttt{SECCOMP} decision, and the
audit trail will faithfully record a false provenance for it --- the log
looks legitimate even though nothing verified the claim.

\textbf{Suggested fix.} Enforce \texttt{AuthorizedBy} against a configured
allowlist of known agent identities before acting on destructive actions,
or use real transport-level authorization (e.g. \texttt{SO\_PEERCRED} UID
checks on the UDS socket) instead of a client-supplied string.

\entrysep

\subsection*{BUG-011 --- No Protected-PID Gate on the Go Control-Plane Side}
\addcontentsline{toc}{subsection}{BUG-011 --- No protected-PID gate (Go side)}
\bugmeta{High (Security)}{kb-control-plane/internal/enforcement/enforce.go (Contain), called from SetContainment and SubmitAgentDecision}{Open}

\textbf{Description.} \texttt{kb-core}'s CPM (Critical Process Module)
correctly gates containment writes on the C sensor side --- refusing to
contain PID 1, kernel threads, or protected executables. But the Go
control plane is the \emph{other} entry point a containment request can
originate from (operator RPC, agent decision), and it has \textbf{zero
exemption logic of its own}.

\textbf{Confirmed in source.} Grepping \texttt{enforce.go} for any
PID-1/self-PID/protected-list check returns nothing --- \texttt{Contain()}
forwards any PID it is given exactly the same way, relying entirely on
\texttt{kb-core}'s gate one hop downstream over \texttt{kbct.sock}.

\textbf{Impact.} A single bug or race in \texttt{kb-core}'s CPM gate has
no backstop on the Go side. A request to contain PID 1, \texttt{systemd},
or \texttt{kbd}'s own PID is forwarded identically to a request against
any other process.

\textbf{Suggested fix.} A minimal check in \texttt{Contain()} (or
immediately before it) refusing/alerting on requests targeting PID 1,
\texttt{os.Getpid()} (the control plane itself), and a small configured
list of critical process names --- independent of whatever \texttt{kb-core}
does downstream.

\entrysep

\subsection*{BUG-012 --- No Input Validation or Rate Limiting on Containment RPCs}
\addcontentsline{toc}{subsection}{BUG-012 --- No input validation/rate limiting}
\bugmeta{Medium (Security/Robustness)}{kb-control-plane/internal/controlplane/grpc.go, SetContainment + SubmitAgentDecision}{Open}

\textbf{Description.} Both RPCs accept any \texttt{int32}/\texttt{uint32}
\texttt{Pid} (including \texttt{0}, a negative cast, or a PID that does
not exist) and act on it without checking existence against the store
first. Neither has any throttle.

\textbf{Confirmed in source:}
\begin{lstlisting}[language=Go]
func (cp *ControlPlane) SetContainment(ctx context.Context,
    req *pb.ContainmentRequest) (*pb.ContainmentResponse, error) {
    cp.enforcer.Contain(req.Pid, uint32(req.Level), req.Reason)
    ...
\end{lstlisting}
No \texttt{store.GetProcessState} existence check precedes the
\texttt{Contain()} call in either RPC; no rate-limiter or token bucket
exists anywhere in this file.

\textbf{Impact.} A buggy or compromised agent client can call
\texttt{SubmitAgentDecision} in a tight loop with no backpressure, and
containment actions can be issued against PIDs the store has never even
observed.

\textbf{Suggested fix.} Validate \texttt{Pid} exists via
\texttt{store.GetProcessState} before acting; add a simple per-agent-ID
token-bucket or minimum-interval check.

\entrysep

\subsection*{BUG-013 --- SSH Listener Has No Auth-Attempt Throttling; Dev-Mode Fallback Silently Weakens Key Material}
\addcontentsline{toc}{subsection}{BUG-013 --- SSH throttling / dev-mode fallback}
\bugmeta{Medium-High (Security)}{kb-control-plane/internal/ssh/server.go, internal/ssh/config.go}{Open (architectural fix planned, not yet implemented)}

\textbf{Description.} This is the daemon's only network-exposed listener
(port 2222) --- every other socket is UDS, local-only. Two distinct
issues:

\begin{enumerate}
\item \textbf{No rate limiting/lockout on failed public-key attempts.}
\texttt{wish.WithPublicKeyAuth} calls \texttt{ValidatePublicKey} per
attempt with no counter, backoff, or fail2ban-style hook. Public-key auth
makes brute-force cryptographically impractical, but there is no
protection against connection-exhaustion/resource-abuse from repeated
handshakes.
\item \textbf{Dev-mode fallback is a real security downgrade, not just a
convenience.} When \texttt{/etc/kb} is not writable and
\texttt{KB\_DEV=true}, \texttt{LoadConfig()} silently switches to host
keys and an \texttt{authorized\_keys} file stored \textbf{relative to the
process's current working directory}. If \texttt{kbd} is ever started
with \texttt{KB\_DEV=true} from an unexpected working directory (a
misconfigured systemd unit, a CI job, a debugging session on a
semi-production box), the host key silently regenerates or resolves
somewhere unexpected --- breaking host-key pinning for anyone who
previously trusted the old fingerprint.
\end{enumerate}

\textbf{Impact.} An attacker with network reach to port 2222 has unlimited
authentication attempts against the SSH listener; a misconfigured
dev-mode deployment can silently rotate host keys in a way that trains
operators to ignore host-key-changed warnings (classic MITM setup).

\textbf{Suggested fix.} Short-term --- per-IP attempt counter/backoff, and
require an explicit second opt-in (e.g. \texttt{KB\_SSH\_ALLOW\_DEV\_FALLBACK=true})
beyond just \texttt{KB\_DEV=true} before falling back to CWD-relative
keys. Longer-term (already decided, not yet built): replace the custom Go
SSH server with real OpenSSH + \texttt{ForceCommand}, inheriting
\textasciitilde 25 years of hardening instead of re-implementing
session/auth handling in a far less battle-tested library.

\entrysep

\subsection*{BUG-014 --- kb-tui Launch Path Is a Hardcoded, Single-Developer Absolute Path}
\addcontentsline{toc}{subsection}{BUG-014 --- Hardcoded kb-tui dev path}
\bugmeta{Medium (Deployment)}{kb-control-plane/internal/ssh/session.go}{Open}

\textbf{Description.} The SSH session handler locates the \texttt{kb-tui}
binary via a hardcoded fallback chain that includes an absolute path
specific to one developer's machine, pointing at a Cargo \textbf{debug}
build.

\textbf{Confirmed in source:}
\begin{lstlisting}[language=Go]
tuiPath := os.Getenv("KB_TUI_PATH")
// falls back to:
"/home/emergence/Desktop/kernel-borderlands/kb-op/kb-tui/target/debug/kb-tui",
"/home/emergence/Desktop/kernel-borderlands/kb-op/kb-tui/kb-tui",
"kb-tui",
\end{lstlisting}

\textbf{Impact.} Deployed to any other host, every SSH session either
silently falls through to a \texttt{\$PATH} lookup (works only if someone
manually installed \texttt{kb-tui} there) or fails outright at SSH-login
time --- there is no build- or install-time contract keeping
\texttt{kb-control-plane} (Go) and \texttt{kb-tui} (Rust) in sync, since
they are two separately-built subsystems glued together only by this
file path.

\textbf{Suggested fix.} Standardize on one fixed install-time location
(e.g. \texttt{/usr/local/bin/kb-tui}), matching how \texttt{kbd} itself
would be installed. Keep only \texttt{KB\_TUI\_PATH} (override) and the
fixed path (common case); drop the dev-machine absolute path and
bare-\texttt{\$PATH} fallbacks entirely. Fail loudly with a specific
message instead of silently trying three guesses.

\entrysep

\subsection*{BUG-015 --- kbctl audit verify Bypasses Its Own Deferred Cleanup on a Detected Chain Break}
\addcontentsline{toc}{subsection}{BUG-015 --- kbctl os.Exit skips cleanup}
\bugmeta{Low}{kb-op/kbctl/cmd\_audit.go, auditVerifyCmd}{Open}

\textbf{Description.} On a detected tamper (\texttt{!resp.ChainIntact}),
the handler calls \texttt{os.Exit(1)} directly:
\begin{lstlisting}[language=Go]
fmt.Printf("CHAIN BROKEN -- %d entries verified before break\n",
    resp.EntriesVerified)
if resp.Error != "" {
    fmt.Println(resp.Error)
}
os.Exit(1)
return nil
\end{lstlisting}
\texttt{os.Exit()} terminates the process immediately without running any
deferred calls --- so the \texttt{defer conn.Close()} and
\texttt{defer cancel()} registered earlier in the same function never
execute.

\textbf{Impact.} Minor in practice (a short-lived CLI process, the OS
reclaims the socket/fd on exit regardless), but it is a real defect and
notably sloppy in exactly the one code path whose entire purpose is
flagging tamper events --- the most security-relevant command in the CLI
has the least-clean shutdown of any command in the file.

\textbf{Suggested fix.} \texttt{return fmt.Errorf(...)} and let
\texttt{main()}'s existing \texttt{os.Exit(1)} (after printing to stderr)
handle the non-zero exit, same as every other command in this CLI already
does --- rather than calling \texttt{os.Exit} directly from inside
\texttt{RunE}.

\clearpage

% ================= RESOLVED =================
\section{Resolved Since the Gaps Were First Documented}

These were open at the time \texttt{control-plane-catalog.md} was written
but are confirmed fixed in current source --- included to show real
progress, not just outstanding issues.

\subsection*{RESOLVED-1 --- Audit Hash-Chain Verification Was Implemented but Never Wired to Anything Callable}
\addcontentsline{toc}{subsection}{RESOLVED-1 --- Audit chain verification wired}
\textbf{Was:} \texttt{Logger.VerifyChain()} existed and was correct, but
had zero callers outside its own test --- tampering with the audit log in
SQLite directly would never be detected or alerted on.

\textbf{Now:} Confirmed wired on both paths --- \texttt{GET
/api/audit/verify} (\seqsplit{\texttt{internal/controlplane/http.go}}) and a
\texttt{VerifyAuditChain} gRPC RPC
(\texttt{internal/controlplane/grpc.go}, \texttt{proto/kb.proto}), plus a
startup check in \texttt{controlplane.go}.

\subsection*{RESOLVED-2 --- GetProcessState Did Not Distinguish "Not Found" From "Zero Value"}
\addcontentsline{toc}{subsection}{RESOLVED-2 --- GetProcessState NotFound}
\textbf{Was:} A PID never tracked by the store returned an empty-but-valid
\texttt{ProcessState} struct indistinguishable from a real all-defaults
process.

\textbf{Now:} Confirmed fixed --- \texttt{GetProcessState} returns
\texttt{status.Errorf(codes.NotFound, "no tracked process with
pid=\%d", req.Pid)} when the store lookup misses.

\subsection*{RESOLVED-3 --- Policy Had No Hot-Reload Path on the Go Side}
\addcontentsline{toc}{subsection}{RESOLVED-3 --- Policy hot-reload}
\textbf{Was:} \texttt{control-plane-catalog.md} \S 2.6 documented that
changing \texttt{policy.yaml} required a full \texttt{kbd} restart --- the
Go-side \texttt{Engine} had no reload path at all.

\textbf{Now:} Confirmed fixed and correctly implemented ---
\texttt{kbctl policy reload} $\rightarrow$ \texttt{ReloadPolicy} gRPC
handler $\rightarrow$ \texttt{reloadPolicyFromDisk()} builds a fresh
\texttt{policy.Engine} and swaps it into \texttt{cp.policy} under a proper
\texttt{sync.RWMutex}. The write side takes a full lock; the read side (in
the zone-transition handler) takes a read lock --- genuinely race-free,
not just apparently working. Found while auditing \texttt{kbctl} for
BUG-015; initially suspected a data race here, traced it fully, and it
checks out.

\clearpage

% ================= HISTORICAL =================
\section{Historical --- Critical Bugs From CHANGELOG.md, Already Fixed and Verified Live}

\subsection*{HIST-1 --- sudo/PAM System-Wide Lockout During Live LSM Enforcement}
\addcontentsline{toc}{subsection}{HIST-1 --- sudo/PAM lockout}
\bugmeta{Critical}{kb-core/ebpf/kbd\_sensor.bpf.c}{Fixed}
\textbf{Root cause:} \texttt{kb\_lsm\_file\_open}'s sensitive-path block
was unconditional --- applied to every process, not just contained ones
--- inconsistent with every other LSM hook in the same file. The first
time it fired correctly, it blocked \texttt{/etc/sudoers}/\texttt{/etc/shadow}
reads system-wide, breaking \texttt{sudo}/\texttt{su}/\texttt{pkexec}/login
PAM simultaneously, with no already-authenticated root path left inside
the machine to recover from --- required a full external VM restart.

\textbf{Fix:} Made the block containment-gated (level $\geq 2$ only),
matching the rest of the file.

\textbf{Verification:} \texttt{sudo whoami} succeeds with the sensor
actively enforcing; a deliberately-contained test process still has its
access correctly blocked.

\subsection*{HIST-2 --- SetContainment Silently Never Applied Containment at All}
\addcontentsline{toc}{subsection}{HIST-2 --- SetContainment silent failure}
\bugmeta{Critical}{kb-control-plane / kb-core wire protocol}{Fixed}
\textbf{Root cause:} \texttt{MsgTypeContainmentCmd} was \texttt{3} on the
Go side but the C sensor checked for \texttt{KB\_WIRE\_MSG\_CONTAINMENT\_CMD}
(\texttt{5}) --- every containment command from \texttt{kb-tui}/\texttt{kbctl}
was silently dropped, while \texttt{kbd} logged a successful audit entry
regardless (the sensor never NACKs), completely masking the failure.

\textbf{Fix:} Corrected the Go-side constant to \texttt{5}.

\textbf{Verification:} A test PID placed into Seccomp containment via
\texttt{kb-tui} had its sensitive-path access correctly blocked
afterward; an uncontained process reading the same file continued to
succeed.

\subsection*{HIST-3 --- Sensitive-Path LSM Block Never Actually Enforced, Despite Being Documented as Active}
\addcontentsline{toc}{subsection}{HIST-3 --- LSM zero-padding bug}
\bugmeta{High}{kb-core/ebpf/kbd\_sensor.bpf.c}{Fixed}
\textbf{Root cause:} \texttt{bpf\_d\_path()} writes a NUL-terminated path
but does not guarantee the buffer's tail past the terminator is zeroed;
the map lookup compared the full fixed-size 64-byte key including that
uninitialized tail, so even an exact, correctly-registered path never
matched.

\textbf{Fix:} Explicitly zero-fill \texttt{path\_buf[len..63]} after
\texttt{bpf\_d\_path()} returns.

\end{document}
